\documentclass[]{article}

%opening
\title{\textbf{LA RETRO-INGENIERIE D’UNE ORDONNANCE DE RENVOI}}
\date{}
\begin{document}
\thispagestyle{empty}
\begin{center}
	\rule{1\textwidth}{2pt}\vspace{\baselineskip}
	\vspace{\baselineskip}\huge\textbf{THEORIES ET PRATIQUES DE L’INVESTIGATION  NUMERIQUE.}
	\rule{1\textwidth}{2pt}\vspace{\baselineskip}
	\vspace{\baselineskip}
	\vfill
	\textbf{Enseignant: M.MINKA THIERRY.}\\
	\vfill
	Par l'étudiante: CHEUTCHEU LONGA HILARY PAULA\\
	Filière CIN 4\vfill
	\vspace{\baselineskip}
	\emph{\textbf{Année Academique: 2025-2026.}}
	\rule{1\textwidth}{2pt}\vspace{\baselineskip}
\end{center}
\maketitle
\section*{INTRODUCTION}
	La rétro-ingénierie est une méthode qui tente d’expliquer par déduction et analyse systémique, comment un mécanisme, un dispositif, un système ou un programme existant, accompli tune tache ; sans connaissance précise de la manière dont il fonction. En investigation numérique elle est définit comme l’analyse d’un évènements en étudiant les différentes étapes de sa réalisation.
	
\subsection*{CONTEXT DU DEVOIR}
	Pour ce devoir, il nous a été demandé de faire du “reverse engineering” d’une ordonnance de renvoi. C’est-à-dire déterminer tous les éléments qui ont été présenté au juges avant l’obtention de cette ordonnance. Après lecture de l’ordonnance mise à notre disposition, la plupart des éléments y sont cités. Ces éléments qui ont permis de reconstituer à peu près exactement les faits. \\

\newpage
\subsection*{ELEMENTS CONSTITUTIFS DE L'ORDONNACE}
\paragraph{}
Pour commencer, l’expert a examiné les preuves qui était à sa disposition tels que les appareils mobiles de la victime pour retracer les derniers conversation, appels, message et autres données qui pouvait être utile et les document d’autopsie de ce dernier pour pouvoir avoir un aperçu de ce qui est arrivé. Les résultats d’analyse des données du téléphone de la victime ont permis d’avoir des suspects et/ou des témoins pour en savoir plus. Delà a suivi les interrogatoires avec les présumés suspects ou témoignes qu’on relevé du mobile de la victime, tous ceci est suivi d’une analyse des différentes données des appareils de ses derniers pour pouvoir déterminer un lien quel qu’on que.
\paragraph{}
Des analyses de géolocalisation des personnes impliquées ainsi que les examens des vidéos de surveillance ont été réalisé pour avoir non seulement des preuves mais aussi pour déterminer le quel des aveux lors des interrogatoire son vérité et permettant donc de reconstituer les faits. Les interviews ont été faits avec ses individus qui devait être entendu pour permettre d’encaisser certains incompréhension et manquement liées à l’affaire. Ensuite sont venu les interrogatoires avec confrontation pour dissocier la version viridité du faux, et l’examens des différents comptes mobiles des suspects pour analyser les différentes transactions effectuer par ces derniers durant ces périodes.

\paragraph{}
Bien que le processus d’exécution soit assez long et nécessite assez de compétences, nous pouvons ressemble les diffèrent éléments présenté comme en étant :\\
i.	Examens des données du téléphone de la victime et de suspects ainsi que des listings et historique auprès des opérateurs mobiles,\\
ii.	Examens des données de géolocalisation,\\
iii.	Interrogatoires et interviews,\\
iv.	Examens des vidéos de surveillances,\\
v.	Examens des comptes bancaire et transactions,\\
vi.	Interrogatoire avec confrontation,\\
vii.	Reconstitution des faits avec preuve et aveu.\\

\newpage
\section*{CONCLUSION}
La lecture du document a permis de découvrir la plupart des outils et techniques utilisé pour avoir un certain nombre d'information compromettant(preuves) qui ont permis de reconstituer les évènements qui ont eu lieu et d'avoir cette ordonnance. Des sanctions ont été attribués à chacun des suspects en fonction de leurs infractions et des sanctions cités dans le code pénal et le code de justice militaire.
\end{document}
